\section{A Body Made in Three Months, and What Was Made}

Between 2 July and 18 October 1988 a group of priests who had just left
the Society of Saint Pius X acquired, from the Holy See, something the
body they had left had not held since 1975: a canonical existence. The
speed is the first thing anyone notices. Twelve or thirteen clerics
signed a foundation act in a Cistercian abbey on 18 July; a Roman
commission answered them four days later; a liturgical decree followed on
10 September; and on 18 October the Priestly Fraternity of Saint Peter
was erected as a clerical society of apostolic life of pontifical right,
with constitutions approved for three years and a superior general named
for the same term.

The speed is also the first thing that misleads. It suggests a
concession hurriedly granted in the aftermath of a rupture, and therefore
a concession that a later pope could as hurriedly withdraw. That reading
has been the common property of the Fraternity's critics on both flanks
for nearly forty years: to the Society of Saint Pius X it proves that the
foundation was a Roman device to divide the traditionalist movement; to
the Fraternity's Roman critics it proves that the whole arrangement was
transitional, a bridge for people shocked in 1988 and not a permanent
feature of the Latin Church. Both readings treat what happened in 1988 as
an act of favour.

This account asks a narrower and more answerable question. \emph{What,
exactly, did the acts of 1988 constitute, and what has kept that
constitution in being through four subsequent changes in the Roman law of
the liturgy?} The question can be answered from documents, because the
governing acts exist and one of them---the act that matters most---was
printed in \work{Acta Apostolicae Sedis}.

Six subordinate questions follow. Under what authority did a pontifical
commission created for reconciliation erect an institute at all? What did
the erection decree of 18 October 1988 actually grant, and in what terms?
What happened in 1999 and 2000, when the Holy See told the Fraternity's
superiors that they could not require of their own members the exclusive
liturgical practice for which the Fraternity had been founded, and then
removed a superior general and appointed his successor by decree? What
did the definitive approval of the constitutions in 2003 settle? What did
\work{Traditionis custodes} do to a body erected by the Pontifical
Commission \work{Ecclesia Dei}, and what did the decree of 11 February
2022 do about it? And what is the Fraternity's position, in law, at this
account's cutoff of 25 July 2026, with an apostolic visitation open and
no published outcome?

\historythesis{The Fraternity of Saint Peter was not given a privilege in
1988; it was given a constitution, and the two have different legal
lives. The papal rescript \latin{Quia peculiare munus} of 18 October
1988---printed in \work{Acta Apostolicae Sedis} 82 (1990) 533--534---gave
the Pontifical Commission \work{Ecclesia Dei} the faculty to erect the
Fraternity as a society of apostolic life of pontifical right ``with
those particular marks of which \work{Ecclesia Dei} n.~6~a speaks,'' and
in the same breath made one thing---and only one thing---expressly
provisional: the Commission's own exercise of the authority of the Holy
See over such bodies, granted \latin{donec aliter provideatur}, until
otherwise provided. Every later Roman act that has moved the Fraternity
has moved that provisional element and not the erection. The decommission
of the Commission in 2019 cites n.~6 of the rescript by number;
\work{Traditionis custodes} art.~6 transfers the institutes erected by
the Commission to another dicastery; the apostolic visitation of 2024 is
opened by that dicastery. The one occasion on which the Holy See reached
past the supervisor and into the substance was 1999--2000, and what it
refused then was not the Fraternity's liturgical patrimony but the
attempt to make that patrimony a personal obligation binding each member
against the universal law. The unresolved question at this account's
cutoff is not whether the Fraternity exists but what the decree of 11
February 2022 is: a papal faculty granted \latin{extra ordinem} and never
promulgated in the official gazette, whose only reference to
\work{Traditionis custodes} is a verb of advice.}

The method is documentary and the hierarchy of witnesses is stated once
here and applied throughout. Acts printed in \work{Acta Apostolicae
Sedis} control what they record. Official texts published by the Holy See
on its own site or in its press bulletin control acts not printed in the
gazette. Where the only available text of a Roman act is one published by
the Fraternity itself, or by a third party sympathetic or hostile, this
account says so at the point of use, names the host, and does not let a
party text establish a contested fact against or in the absence of an
official one. Documents published by the Fraternity establish the
Fraternity's own positions, self-description, internal events and
self-reported numbers, and nothing else.

Three warnings about vocabulary belong at the front. ``Society of
apostolic life'' is a canonical classification with a defined content in
canons 731--746 of the 1983 Code; it is not a religious order, its members
take no religious vows, and nothing in this account should be read as
saying otherwise. ``Pontifical right'' means that the Apostolic See
erected the body and that its internal government and discipline are
immediately and exclusively subject to that See under canon 593; it does
not mean exemption from the diocesan bishop in public worship, the care
of souls, or apostolic works. And ``the 1962 books'' means the Missal,
Ritual, Pontifical and Roman Breviary in force in that year, named as
such in the acts that grant them; ``the traditional Latin Mass'' is a
looser expression that this account avoids wherever the identity of a
particular book or faculty carries weight.

Finally, a boundary. The history of the Society of Saint Pius X is the
subject of the companion volume in this series and is not retold here.
This account touches that history at exactly one point---the branch of
1988---and again, glancingly, at the events of July 2026, where the
question is only what the Fraternity did and did not say. No judgment is
offered here on the canonical position of the Society of Saint Pius X;
where the Fraternity's own founders published such a judgment, it is
reported as theirs.
